\section{Limitations, Open Questions \& Conclusion}

\paragraph{Base-case limitation.}
The theorem is an existence result in the Sen24 base case only.
It does not claim that raw repair explanation is non-invariant for every encoding comparison, every theorem family, or every finite impossibility witness.
What is established is that the phenomenon already occurs in a tightly controlled base-case artifact bundle.

\paragraph{Local-rationality limitation.}
All arguments in this paper remain inside the current local-rationality setting induced by \code{no\_cycle3} and \code{no\_cycle4}.
Those levers are local forbidden-pattern approximations, not full \code{SocialAcyclic}.
Accordingly, the paper makes no full-acylcity claim and no full Sen-family or Arrow-family generality claim.

\paragraph{What remains open.}
The main open problem is the repair-canonicalization layer itself.
This paper distinguishes raw repair explanation from grouped repair explanation, but it does not characterize when grouped repair explanations are invariant.
At most, future work may identify sufficient conditions under explicit grouping contracts.
That is an open question, not a theorem claimed here.

\paragraph{Connection to future repair geometry work.}
The next conceptual step is not necessarily another base-case witness.
It is to study the geometry induced by abstraction contracts:
which contracts preserve diagnosability, which collapse substantively different repairs, and which align with the social-choice semantics that a paper or audit report wants to foreground.
That future work may connect to richer repair geometry, but it should build on the present distinction rather than erase it.

\paragraph{Conclusion.}
The Candidate B witness shows that clause-equivalent bundled and split encodings can agree on the existence of an impossibility boundary while disagreeing on the raw lever-level minimal repair family.
That is enough to separate finite witness legitimization from repair presentation canonicality in the current repository evidence base.
The main implication is normative: one may have a trustworthy finite witness without yet having a canonical repair explanation.

Proof-carrying or witness-carrying finite impossibility evidence can legitimately certify that a boundary is real.
It does not, by itself, determine how repairs should be grouped, named, or presented to a reviewer, designer, or downstream system.
For audit-oriented repair presentation, an explicit abstraction contract is therefore not optional gloss but part of the reporting specification itself.
